\documentclass[11pt,a4paper]{article}
\usepackage[utf8]{inputenc}
\usepackage[T1]{fontenc}
\usepackage[portuguese]{babel}
\usepackage{geometry}
\geometry{margin=2.5cm}
\usepackage{listings}
\usepackage{xcolor}
\usepackage{amsmath}
\usepackage{amssymb}
\usepackage{hyperref}
\usepackage{enumitem}
\usepackage{tcolorbox}
\usepackage{tabularx}
\usepackage{booktabs}

\definecolor{codebg}{rgb}{0.95,0.95,0.95}
\definecolor{codegreen}{rgb}{0.1,0.5,0.1}
\definecolor{codepurple}{rgb}{0.5,0.1,0.5}
\definecolor{codegray}{rgb}{0.5,0.5,0.5}

\lstdefinestyle{python}{
    backgroundcolor=\color{codebg},
    commentstyle=\color{codegray}\itshape,
    keywordstyle=\color{codepurple}\bfseries,
    stringstyle=\color{codegreen},
    basicstyle=\ttfamily\small,
    breaklines=true,
    frame=single,
    language=Python,
    showstringspaces=false,
    tabsize=4,
    literate={á}{{\'a}}1 {ã}{{\~a}}1 {é}{{\'e}}1 {í}{{\'i}}1 {ó}{{\'o}}1 {õ}{{\~o}}1 {ú}{{\'u}}1 {ç}{{\c{c}}}1 {Á}{{\'A}}1 {Ã}{{\~A}}1 {É}{{\'E}}1 {Í}{{\'I}}1 {Ó}{{\'O}}1 {Õ}{{\~O}}1 {Ú}{{\'U}}1 {Ç}{{\c{C}}}1 {â}{{\^a}}1 {ê}{{\^e}}1 {ô}{{\^o}}1 {û}{{\^u}}1 {Â}{{\^A}}1 {Ê}{{\^E}}1 {Ô}{{\^O}}1 {Û}{{\^U}}1 {à}{{\`a}}1 {è}{{\`e}}1 {ì}{{\`i}}1 {ò}{{\`o}}1 {ù}{{\`u}}1
}
\lstset{style=python}

\newtcolorbox{objetivos}[1][]{
  colback=green!5!white,
  colframe=green!40!black,
  title=O que você vai aprender nesta página,
  fonttitle=\bfseries,
  #1
}

\newtcolorbox{exercicio}[1]{
  colback=blue!5!white,
  colframe=blue!40!black,
  title=#1,
  fonttitle=\bfseries
}

\newtcolorbox{solucao}{
  colback=yellow!5!white,
  colframe=orange!60!black,
  title=Solução proposta,
  fonttitle=\bfseries
}

\title{\textbf{Data Analysis with Python 2024--2025}\\Parte 3: NumPy (Parte 1)\\\large Conteúdo Teórico}
\author{Tradução e Adaptação: Universidade de Helsinque (MOOC.fi)}
\date{}

\begin{document}

\maketitle

\section*{Nota Importante}
Este material é uma tradução e adaptação baseada no curso \textit{Data Analysis with Python 2024-2025}, oferecido pela \textbf{Universidade de Helsinque (MOOC.fi)}.

\begin{objetivos}
\begin{itemize}
\item Saber o que é o NumPy.
\item Entender por que o NumPy é eficiente.
\item Saber criar, remodelar, acessar, combinar, dividir e agregar \emph{arrays}.
\end{itemize}
\end{objetivos}

\tableofcontents
\newpage

\section{Introdução}

O NumPy é uma biblioteca Python para trabalhar com \emph{arrays}
multidimensionais. Ele contém tanto as estruturas de dados necessárias
para armazenar e acessar esses arrays quanto as operações e funções para
realizar cálculos usando-os. Embora os arrays sejam normalmente usados
para armazenar números, outros tipos de dados também podem ser
armazenados, como strings. Diferentemente das listas do Python puro, a
estrutura de dados fundamental do NumPy, o \emph{array}, precisa ter o
mesmo tipo de dado em todos os seus elementos. Essa homogeneidade
permite a criação de funções altamente otimizadas, que recebem e
retornam arrays.

Há vários usos para arrays de alta dimensão em análise de dados. Por
exemplo, eles podem ser usados para:

\begin{itemize}
\item armazenar matrizes, resolver sistemas de equações lineares,
  encontrar autovalores/autovetores, encontrar decomposições de matrizes
  e resolver outros problemas conhecidos da álgebra linear;
\item armazenar dados de medições multidimensionais. Por exemplo, um
  elemento \texttt{a[i,j]} em um array bidimensional pode armazenar a
  temperatura $t_{ij}$ medida nas coordenadas $i, j$ de uma superfície
  bidimensional;
\item representar imagens e vídeos:
  \begin{itemize}
  \item uma imagem em escala de cinza pode ser representada como um
    array bidimensional;
  \item uma imagem colorida pode ser representada como um array
    tridimensional, em que a terceira dimensão contém os componentes de
    cor vermelho, verde e azul;
  \item um vídeo colorido pode ser representado como um array
    quadridimensional;
  \end{itemize}
\item uma tabela bidimensional pode armazenar uma sequência de amostras,
  cada uma dividida em características (\emph{features}). Por exemplo,
  poderíamos medir as condições climáticas uma vez por dia, incluindo
  temperatura, direção e velocidade do vento, e quantidade de chuva.
  Teríamos, então, uma amostra por dia, e as características seriam
  temperatura, vento e chuva. Na representação padrão desse tipo de
  dado tabular, as linhas correspondem às amostras e as colunas às
  características. Veremos mais desse tipo de dado nos capítulos sobre
  Pandas e Scikit-learn.
\end{itemize}

Neste capítulo, veremos:

\begin{itemize}
\item criação de arrays;
\item tipos e atributos de arrays;
\item acesso a arrays com indexação e \emph{slicing};
\item remodelagem (\emph{reshaping}) de arrays;
\item combinação e divisão de arrays;
\item operações rápidas em arrays;
\item agregações de arrays;
\item regras das operações binárias entre arrays;
\item operações matriciais conhecidas da álgebra linear.
\end{itemize}

Começamos importando a biblioteca NumPy, usando a abreviação padrão
\texttt{np}:

\begin{lstlisting}
import numpy as np
\end{lstlisting}

\section{Criação de arrays}

Existem várias formas de criar arrays NumPy. Uma delas é passar uma
lista (possivelmente aninhada) como parâmetro ao construtor
\texttt{array}:

\begin{lstlisting}
np.array([1,2,3])   # array unidimensional
# array([1, 2, 3])
\end{lstlisting}

Note que omitir os colchetes na expressão acima, ou seja, chamar
\texttt{np.array(1,2,3)}, resulta em erro.

Um array bidimensional pode ser criado listando as linhas do array:

\begin{lstlisting}
np.array([[1,2,3], [4,5,6]])
# array([[1, 2, 3],
#        [4, 5, 6]])
\end{lstlisting}

Da mesma forma, um array tridimensional pode ser descrito como uma
lista de listas de listas:

\begin{lstlisting}
np.array([[[1,2], [3,4]], [[5,6], [7,8]]])
# array([[[1, 2],
#         [3, 4]],
#
#        [[5, 6],
#         [7, 8]]])
\end{lstlisting}

Existem também algumas funções auxiliares para criar tipos comuns de
arrays:

\begin{lstlisting}
np.zeros((3,4))
# array([[ 0.,  0.,  0.,  0.],
#        [ 0.,  0.,  0.,  0.],
#        [ 0.,  0.,  0.,  0.]])
\end{lstlisting}

Para especificar que os elementos são \texttt{int}, em vez de
\texttt{float}, usa-se o parâmetro \texttt{dtype}:

\begin{lstlisting}
np.zeros((3,4), dtype=int)
# array([[0, 0, 0, 0],
#        [0, 0, 0, 0],
#        [0, 0, 0, 0]])
\end{lstlisting}

De forma semelhante, \texttt{ones} inicializa todos os elementos com o
valor um, \texttt{full} inicializa todos os elementos com um valor
especificado, e \texttt{empty} deixa os elementos sem inicialização:

\begin{lstlisting}
np.ones((2,3))
# array([[ 1.,  1.,  1.],
#        [ 1.,  1.,  1.]])

np.full((2,3), fill_value=7)
# array([[7, 7, 7],
#        [7, 7, 7]])

np.empty((2,4))
# valores nao inicializados (lixo de memoria)
\end{lstlisting}

A função \texttt{eye} cria a matriz identidade, ou seja, uma matriz cujos
elementos na diagonal valem um, e os demais valem zero:

\begin{lstlisting}
np.eye(5, dtype=int)
# array([[1, 0, 0, 0, 0],
#        [0, 1, 0, 0, 0],
#        [0, 0, 1, 0, 0],
#        [0, 0, 0, 1, 0],
#        [0, 0, 0, 0, 1]])
\end{lstlisting}

A função \texttt{arange} funciona como a função \texttt{range}, mas
produz um array em vez de uma lista:

\begin{lstlisting}
np.arange(0,10,2)
# array([0, 2, 4, 6, 8])
\end{lstlisting}

Para intervalos com números não inteiros, é preferível usar
\texttt{linspace}:

\begin{lstlisting}
np.linspace(0, np.pi, 5)   # intervalo com 5 elementos igualmente espacados
# array([ 0. , 0.78539816, 1.57079633, 2.35619449, 3.14159265])
\end{lstlisting}

Com \texttt{linspace} não é preciso calcular o tamanho do passo: em vez
disso, especifica-se a quantidade desejada de elementos. Por padrão, o
ponto final é incluído no resultado -- diferentemente do que ocorre com
\texttt{arange}.

\subsection{Arrays com elementos aleatórios}

Para testar nossos programas, podemos usar dados reais como entrada.
Porém, dados reais nem sempre estão disponíveis, e coletá-los pode levar
tempo. Podemos, em vez disso, gerar números aleatórios como substituto.
Eles podem ser gerados facilmente com o NumPy, a partir de diversas
distribuições diferentes -- mencionamos apenas algumas a seguir. Dados
aleatórios podem simular dados reais melhor do que, por exemplo,
intervalos regulares ou arrays de valores constantes. Às vezes também
precisamos de números aleatórios para escolher um subconjunto de dados
reais (amostragem). O NumPy pode produzir facilmente arrays com o
formato desejado preenchidos com números aleatórios. Alguns exemplos a
seguir.

\begin{lstlisting}
np.random.random((3,4))
# elementos distribuidos uniformemente no intervalo semiaberto [0.0, 1.0)

np.random.normal(0, 1, (3,4))
# elementos com distribuicao normal, media 0 e desvio padrao 1

np.random.randint(-2, 10, (3,4))
# elementos inteiros distribuidos uniformemente no intervalo [-2, 10)
\end{lstlisting}

Às vezes é útil poder recriar exatamente os mesmos dados a cada execução
do programa. Por exemplo, se há um erro no programa que só se manifesta
com determinada entrada, depurar o programa exige que ele se comporte de
forma determinística. Podemos gerar números aleatórios de forma
determinística se sempre começarmos a partir do mesmo ponto de partida.
Esse ponto de partida costuma ser um número inteiro, chamado de
\emph{seed} (semente). Exemplo de uso:

\begin{lstlisting}
np.random.seed(0)
print(np.random.randint(0, 100, 10))
print(np.random.normal(0, 1, 10))
\end{lstlisting}

Se você executar essa célula várias vezes, ela sempre produzirá os
mesmos números -- diferentemente dos exemplos anteriores. A chamada a
\texttt{np.random.seed} inicializa o gerador de números aleatórios
\emph{global}. As chamadas \texttt{np.random.random},
\texttt{np.random.normal} etc. usam todas esse gerador global. Também é
possível criar novos geradores de números aleatórios, e usá-los para
amostrar números de uma distribuição:

\begin{lstlisting}
novo_gerador = np.random.RandomState(seed=123)   # RandomState e uma classe;
                                                   # passamos a seed ao construtor
novo_gerador.randint(0, 100, 10)
\end{lstlisting}

Esses recursos serão usados mais adiante no material e nos exercícios,
justamente para que possamos combinar quais são os dados de entrada
aleatórios. Afinal, como poderíamos concordar se um resultado está
correto ou não, se não pudermos concordar sobre qual é a entrada?

\section{Tipos e atributos de arrays}

Um array possui vários atributos: \texttt{ndim} informa o número de
dimensões, \texttt{shape} informa o tamanho em cada dimensão,
\texttt{size} informa o número total de elementos, e \texttt{dtype}
informa o tipo dos elementos. Vamos criar uma função auxiliar para
explorar esses atributos:

\begin{lstlisting}
def info(nome, a):
    print(f"{nome} tem dim {a.ndim}, shape {a.shape}, size "
          f"{a.size}, e dtype {a.dtype}:")
    print(a)

b = np.array([[1,2,3], [4,5,6]])
info("b", b)
# b tem dim 2, shape (2, 3), size 6, e dtype int64:
# [[1 2 3]
#  [4 5 6]]

c = np.array([b, b])   # cria um array tridimensional
info("c", c)
# c tem dim 3, shape (2, 2, 3), size 12, e dtype int64:
# [[[1 2 3]
#   [4 5 6]]
#
#  [[1 2 3]
#   [4 5 6]]]

d = np.array([[1,2,3,4]])   # um vetor-linha
info("d", d)
# d tem dim 2, shape (1, 4), size 4, e dtype int64:
# [[1 2 3 4]]
\end{lstlisting}

Observe como o Python imprimiu o array tridimensional acima. As regras
gerais para imprimir um array de $n$ dimensões como uma lista aninhada
são:

\begin{itemize}
\item a última dimensão é impressa da esquerda para a direita;
\item a penúltima dimensão é impressa de cima para baixo;
\item as demais também são impressas de cima para baixo, e cada
  ``fatia'' é separada da próxima por uma linha em branco.
\end{itemize}

\section{Indexação, \emph{slicing} e remodelagem}

\subsection{Indexação}

Um array unidimensional se comporta como uma lista do Python:

\begin{lstlisting}
a = np.array([1,4,2,7,9,5])
print(a[1])    # 4
print(a[-2])   # 9
\end{lstlisting}

Para um array multidimensional, o índice é uma tupla separada por
vírgulas, em vez de um único inteiro:

\begin{lstlisting}
b = np.array([[1,2,3], [4,5,6]])
print(b)
print(b[1,2])    # linha de indice 1, coluna de indice 2 -> 6
print(b[0,-1])   # linha de indice 0, coluna de indice -1 -> 3

# assim como nas listas, a modificacao por indexacao e possivel
b[0,0] = 10
print(b)
# [[10  2  3]
#  [ 4  5  6]]
\end{lstlisting}

Se fornecermos apenas um índice a um array multidimensional, esse
índice se refere à primeira dimensão do array, ou seja, às linhas:

\begin{lstlisting}
print(b[0])   # primeira linha
print(b[1])   # segunda linha
\end{lstlisting}

\subsection{\emph{Slicing}}

O \emph{slicing} funciona de forma parecida com as listas, mas agora
podemos ter fatias em diferentes dimensões:

\begin{lstlisting}
print(a)          # [1 4 2 7 9 5]
print(a[1:3])     # [4 2]
print(a[::-1])    # inverte o array: [5 9 7 2 4 1]

print(b)
print(b[:,0])     # [10  4]
print(b[0,:])     # [10  2  3]
print(b[:,1:])    # [[2 3]
                   #  [5 6]]
\end{lstlisting}

É possível até atribuir a uma fatia:

\begin{lstlisting}
b[:,1:] = 7
print(b)
# [[10  7  7]
#  [ 4  7  7]]
\end{lstlisting}

Um padrão comum é extrair linhas ou colunas de um array:

\begin{lstlisting}
print(b[:,0])   # primeira coluna
print(b[1,:])   # segunda linha
\end{lstlisting}

\subsection{Remodelagem (\emph{reshaping})}

Quando um array é remodelado, o número de elementos permanece o mesmo,
mas eles são reinterpretados em um novo formato. Um exemplo disso é
interpretar um array unidimensional como um array bidimensional:

\begin{lstlisting}
a = np.arange(9)
anovo = a.reshape(3,3)
info("anovo", anovo)
# anovo tem dim 2, shape (3, 3), size 9, e dtype int64:
# [[0 1 2]
#  [3 4 5]
#  [6 7 8]]
info("a", a)
# a tem dim 1, shape (9,), size 9, e dtype int64:
# [0 1 2 3 4 5 6 7 8]

d = np.arange(4)          # array 1d
dl = d.reshape(1,4)       # vetor-linha
dc = d.reshape(4,1)       # vetor-coluna
info("d", d)
info("dl", dl)
info("dc", dc)
\end{lstlisting}

Uma sintaxe alternativa para criar, por exemplo, vetores-coluna ou
vetores-linha é usar a palavra-chave \texttt{np.newaxis}. Às vezes essa
forma é mais simples ou mais natural do que usar o método
\texttt{reshape}:

\begin{lstlisting}
info("d", d)
info("dlinha", d[np.newaxis, :])
info("dcoluna", d[:, np.newaxis])
\end{lstlisting}

\section{Concatenação, divisão e empilhamento de arrays}

Existem duas formas de combinar vários arrays em um array maior:
\texttt{concatenate} e \texttt{stack}. \texttt{concatenate} recebe
arrays de $n$ dimensões e retorna um array de $n$ dimensões, enquanto
\texttt{stack} recebe arrays de $n$ dimensões e retorna um array de
$n+1$ dimensões. Alguns exemplos:

\begin{lstlisting}
a = np.arange(2)
b = np.arange(2,5)
print(f"a tem shape {a.shape}: {a}")
print(f"b tem shape {b.shape}: {b}")
np.concatenate((a,b))   # concatenando arrays 1d
# array([0, 1, 2, 3, 4])

c = np.arange(1,5).reshape(2,2)
print(f"c tem shape {c.shape}:", c, sep="\n")
np.concatenate((c,c))   # concatenando arrays 2d
# array([[1, 2],
#        [3, 4],
#        [1, 2],
#        [3, 4]])
\end{lstlisting}

Por padrão, \texttt{concatenate} junta os arrays ao longo do eixo 0.
Para juntá-los horizontalmente, adicione o parâmetro \texttt{axis=1}:

\begin{lstlisting}
np.concatenate((c,c), axis=1)
# array([[1, 2, 1, 2],
#        [3, 4, 3, 4]])
\end{lstlisting}

Se quisermos concatenar arrays com dimensões diferentes -- por exemplo,
para adicionar uma nova coluna a um array 2d --, é preciso primeiro
remodelar os arrays para que tenham o mesmo número de dimensões:

\begin{lstlisting}
print("Nova linha:")
print(np.concatenate((c, a.reshape(1,2))))
print("Nova coluna:")
print(np.concatenate((c, a.reshape(2,1)), axis=1))
\end{lstlisting}

Use \texttt{stack} para criar arrays de dimensão maior a partir de
arrays de dimensão menor:

\begin{lstlisting}
np.stack((b,b))
# array([[2, 3, 4],
#        [2, 3, 4]])

np.stack((b,b), axis=1)
# array([[2, 2],
#        [3, 3],
#        [4, 4]])
\end{lstlisting}

A operação inversa de \texttt{concatenate} é \texttt{split}. Seu
argumento especifica o número de partes iguais em que o array será
dividido, ou então especifica explicitamente os pontos de corte.

\begin{lstlisting}
d = np.arange(12).reshape(6,2)
print("d:")
print(d)
d1, d2 = np.split(d, 2)
print("d1:")
print(d1)
print("d2:")
print(d2)

d = np.arange(12).reshape(2,6)
print("d:")
print(d)
partes = np.split(d, (2,3,5), axis=1)
for i, p in enumerate(partes):
    print("parte %i:" % i)
    print(p)
\end{lstlisting}

\section{Computação rápida com funções universais}

Além de fornecer uma forma de armazenar e acessar arrays
multidimensionais, o NumPy também oferece diversas rotinas para realizar
cálculos sobre eles. Uma das razões da popularidade do NumPy é que esses
cálculos podem ser muito eficientes -- muito mais eficientes do que o
que o Python normalmente consegue fazer. Os maiores gargalos de
eficiência são os laços, que podem ser iterados milhões, bilhões, ou até
mais vezes, e por isso devem ser o mais eficientes possível. O que torna
os laços lentos em Python é o fato de ser uma linguagem dinamicamente
tipada: a cada expressão, o Python precisa descobrir os tipos dos
argumentos das operações. Considere o seguinte laço:

\begin{lstlisting}
L = [1, 5.2, "ab"]
L2 = []
for x in L:
    L2.append(x*2)
print(L2)
# [2, 10.4, 'abab']
\end{lstlisting}

A cada iteração desse laço, o Python precisa descobrir o tipo da
variável \texttt{x} -- que, neste exemplo, pode ser um \texttt{int}, um
\texttt{float} ou uma string -- e, dependendo desse tipo, chamar uma
função diferente para realizar a ``multiplicação'' por dois. O que torna
o NumPy eficiente é a exigência de que cada elemento em um array seja do
mesmo tipo. Essa homogeneidade permite criar \emph{operações
vetorizadas}, que não operam sobre elementos individuais, mas sobre
arrays (ou subarrays) inteiros. O exemplo anterior, usando operações
vetorizadas do NumPy:

\begin{lstlisting}
a = np.array([2.1, 5.0, 17.2])
a2 = a*2
print(a2)
# [ 4.2 10.  34.4]
\end{lstlisting}

Como cada iteração usa operações idênticas -- só os dados mudam --, isso
pode ser compilado para código de máquina e executado de uma só vez,
evitando assim a tipagem dinâmica do Python. O nome ``operação vetorial''
vem da álgebra linear, em que a soma de dois vetores $v = (v_1, v_2,
\ldots, v_d)$ e $w = (w_1, w_2, \ldots, w_d)$ é definida elemento a
elemento como $v + w = (v_1+w_1, v_2+w_2, \ldots, v_d+w_d)$.

Além da soma, várias outras funções matemáticas estão definidas em sua
forma vetorial. As operações aritméticas básicas são:

\begin{center}
\begin{tabular}{ll}
\toprule
\textbf{Operação} & \textbf{Operador} \\
\midrule
adição & \texttt{+} \\
subtração & \texttt{-} \\
negação & \texttt{-} \\
multiplicação & \texttt{*} \\
divisão & \texttt{/} \\
divisão inteira & \texttt{//} \\
exponenciação & \texttt{**} \\
resto & \texttt{\%} \\
\bottomrule
\end{tabular}
\end{center}

Elas podem ser combinadas em expressões mais complexas. Um exemplo:

\begin{lstlisting}
b = np.array([-1, 3.2, 2.4])
print(-a**2 * b)
# [   4.41  -80.   -710.016]
\end{lstlisting}

Diversas outras funções matemáticas também estão definidas. Alguns
exemplos:

\begin{lstlisting}
print(np.abs(b))
print(np.cos(b))
print(np.exp(b))
print(np.log2(np.abs(b)))
\end{lstlisting}

Na nomenclatura do NumPy, essas operações vetoriais são chamadas de
\emph{ufuncs} (\emph{universal functions}, funções universais).

\section{Agregações: máximo, mínimo, soma, média, desvio padrão e mais}

Agregações nos permitem condensar as informações de um array em apenas
alguns poucos números.

\begin{lstlisting}
np.random.seed(0)
a = np.random.randint(-100, 100, (4,5))
print(a)
print(f"Minimo: {a.min()}, maximo: {a.max()}")
print(f"Soma: {a.sum()}")
print(f"Media: {a.mean()}, desvio padrao: {a.std()}")
\end{lstlisting}

Em vez de agregar sobre o array inteiro, também podemos agregar apenas
sobre determinados eixos:

\begin{lstlisting}
np.random.seed(9)
b = np.random.randint(0, 10, (3,4))
print(b)
print("Somas das colunas:", b.sum(axis=0))
print("Somas das linhas:", b.sum(axis=1))
\end{lstlisting}

A tabela a seguir resume algumas das funções de agregação mais comuns,
com a função equivalente do Python puro (quando existe), a função do
NumPy e o método correspondente do array:

\begin{center}
\begin{tabular}{lll}
\toprule
\textbf{Função Python} & \textbf{Função NumPy} & \textbf{Método NumPy} \\
\midrule
\texttt{sum} & \texttt{np.sum} & \texttt{a.sum} \\
-- & \texttt{np.prod} & \texttt{a.prod} \\
-- & \texttt{np.mean} & \texttt{a.mean} \\
-- & \texttt{np.std} & \texttt{a.std} \\
-- & \texttt{np.var} & \texttt{a.var} \\
\texttt{min} & \texttt{np.min} & \texttt{a.min} \\
\texttt{max} & \texttt{np.max} & \texttt{a.max} \\
-- & \texttt{np.argmin} & \texttt{a.argmin} \\
-- & \texttt{np.argmax} & \texttt{a.argmax} \\
-- & \texttt{np.median} & -- \\
-- & \texttt{np.percentile} & -- \\
\texttt{any} & \texttt{np.any} & \texttt{a.any} \\
\texttt{all} & \texttt{np.all} & \texttt{a.all} \\
\bottomrule
\end{tabular}
\end{center}

Vamos medir o quanto a função \texttt{sum} do Python é mais lenta do que
a equivalente do NumPy ao agregar um array:

\begin{lstlisting}
a = np.arange(1000)
%timeit np.sum(a)
# 2.37 us +- 33.6 ns por loop (media +- desvio padrao de 7 execucoes,
# 100000 loops cada)

%timeit sum(a)
# 66.6 us +- 792 ns por loop (media +- desvio padrao de 7 execucoes,
# 10000 loops cada)
\end{lstlisting}

A velocidade do NumPy se deve, em parte, ao fato de que seus arrays
precisam ter o mesmo tipo em todos os elementos. Essa exigência permite
diversas otimizações eficientes.

\section{\emph{Broadcasting}}

Já vimos que o NumPy permite operações realizadas elemento a elemento.
Mas o NumPy também permite operações binárias que não exigem que os dois
arrays tenham o mesmo formato. Por exemplo, podemos somar 4 a todos os
elementos de um array com a expressão a seguir:

\begin{lstlisting}
np.arange(3) + np.array([4])
# array([4, 5, 6])
\end{lstlisting}

Na verdade, como um array com apenas um elemento (digamos, 4) pode ser
visto como um escalar 4, o NumPy também aceita a expressão a seguir, que
é equivalente à anterior:

\begin{lstlisting}
np.arange(3) + 4
# array([4, 5, 6])
\end{lstlisting}

Para ter uma ideia de quais operações são permitidas -- ou seja, quais
formatos de dois arrays são \emph{compatíveis} --, é útil pensar que,
antes de realizar a operação binária, o NumPy tenta ``esticar'' os
arrays para que tenham o mesmo formato. No exemplo acima, o NumPy
primeiro esticou o array \texttt{np.array([4])} (ou o escalar 4) para o
array \texttt{np.array([4,4,4])}, e só então realizou a soma elemento a
elemento. No NumPy, esse ``esticamento'' é chamado de
\emph{broadcasting}.

Os arrays de entrada podem, obviamente, ter dimensões maiores, como
mostra o próximo exemplo:

\begin{lstlisting}
a = np.full((3,3), 5)
b = np.arange(3)
print("a:", a, sep="\n")
print("b:", b)
print("a+b:", a+b, sep="\n")
# a:
# [[5 5 5]
#  [5 5 5]
#  [5 5 5]]
# b: [0 1 2]
# a+b:
# [[5 6 7]
#  [5 6 7]
#  [5 6 7]]
\end{lstlisting}

Nesse exemplo, o segundo argumento foi primeiro transformado
(\emph{broadcasted}) no array

\begin{lstlisting}
np.array([[0, 1, 2],
          [0, 1, 2],
          [0, 1, 2]])
\end{lstlisting}

e só então a soma foi realizada. Também pode ser que ambos os arrays de
entrada precisem passar por \emph{broadcasting}, como no exemplo a
seguir:

\begin{lstlisting}
a = np.arange(3)
b = np.arange(3).reshape((3,1))
info("a", a)
info("b", b)
info("a+b", a+b)
# a+b tem dim 2, shape (3, 3), size 9, e dtype int64:
# [[0 1 2]
#  [1 2 3]
#  [2 3 4]]
\end{lstlisting}

Para ver a que formato os argumentos foram convertidos antes da
operação binária, pode-se usar a função \texttt{np.broadcast\_arrays}:

\begin{lstlisting}
a_convertido, b_convertido = np.broadcast_arrays(a,b)
info("a_convertido", a_convertido)
info("b_convertido", b_convertido)
\end{lstlisting}

Assim, os dois arrays passaram por \emph{broadcasting}, mas de formas
diferentes. A seguir, apresentamos as regras que governam o
\emph{broadcasting}:

\begin{enumerate}
\item todos os arrays de entrada com \texttt{ndim} menor que o
  \texttt{ndim} do array de maior dimensão recebem 1's à esquerda de sua
  forma (\emph{shape});
\item o tamanho de cada dimensão do formato de saída é o máximo entre os
  tamanhos de todas as entradas naquela dimensão;
\item uma entrada pode ser usada no cálculo se seu tamanho em uma
  determinada dimensão for igual ao tamanho de saída naquela dimensão,
  ou se for exatamente 1;
\item se uma entrada tiver tamanho 1 em uma dimensão de sua forma, o
  primeiro (e único) valor daquela dimensão é usado em todos os cálculos
  ao longo dela. Em outras palavras, o mecanismo de avanço da
  \emph{ufunc} simplesmente não avança ao longo dessa dimensão (o passo,
  ou \emph{stride}, é 0 nessa dimensão).
\end{enumerate}

Por fim, um exemplo de situação em que os dois arrays não são
compatíveis:

\begin{lstlisting}
a = np.array([1,2,3])
b = np.array([4,5])
try:
    a+b   # isso nao funciona, pois viola a regra 3 acima
except ValueError as e:
    import sys
    print(e, file=sys.stderr)
# operands could not be broadcast together with shapes (3,) (2,)
\end{lstlisting}

\section{Resumo (semana 2)}

\begin{itemize}
\item Aprendemos como expressões regulares podem ser usadas para
  especificar conjuntos regulares de strings:
  \begin{itemize}
  \item sabemos verificar se uma string casa com uma expressão regular;
  \item sabemos extrair trechos de texto que casam com uma expressão
    regular;
  \item sabemos substituir trechos que casam com uma expressão regular
    por outra string.
  \end{itemize}
\item Sabemos ler (e escrever) um arquivo de texto, tanto linha a linha
  quanto o arquivo inteiro de uma vez.
\item Também sabemos especificar a codificação do arquivo. A codificação
  UTF-8 é muito comum e consegue representar quase todo caractere ou
  símbolo de qualquer língua natural.
\item Os parâmetros de um programa ficam no array \texttt{sys.argv}, e
  podemos retornar um valor do programa com a função \texttt{sys.exit}.
\item Os fluxos de arquivo \texttt{sys.stdin}, \texttt{sys.stdout} e
  \texttt{sys.stderr} permitem entrada e saída textual básica de e para
  o programa.
\item Todo valor em Python é um objeto:
  \begin{itemize}
  \item classes são tipos de dados definidos pelo usuário, que
    determinam como construir instâncias (objetos) desse tipo;
  \item a relação entre objetos e classes é testada com
    \texttt{isinstance};
  \item a relação entre classes é testada com \texttt{issubclass}.
  \end{itemize}
\item Exceções sinalizam situações excepcionais, não necessariamente
  erros -- sabemos como capturar e lançar exceções.
\item A eficiência do NumPy se baseia no fato de que as mesmas operações
  podem ser realizadas sobre os elementos de forma rápida, se todos os
  elementos forem do mesmo tipo. São as chamadas operações vetorizadas.
\item Sabemos criar, remodelar, acessar, combinar, dividir e agregar
  arrays.
\end{itemize}

\section*{Referências e créditos}
\addcontentsline{toc}{section}{Referências e créditos}

Este material é uma tradução e adaptação, para a língua portuguesa, do
conteúdo do curso \emph{Data Analysis with Python 2024--2025}, Capítulo
2 -- ``NumPy (Part 1)'', originalmente disponibilizado em inglês pela
plataforma MOOC.fi:

\begin{quote}
\url{https://courses.mooc.fi/org/uh-cs/courses/data-analysis-with-python-2024-2025/chapter-2/numpy}
\end{quote}

\end{document}
